\documentclass[12pt,a4paper]{article}

\usepackage[T1]{fontenc}
\usepackage[utf8]{inputenc}
\usepackage{lmodern}
\usepackage{geometry}
\usepackage{amsmath}
\usepackage{hyperref}
\usepackage{enumitem}
\usepackage{verbatim}

\geometry{margin=2.5cm}

\title{API Mínima da Estrutura de Dados Diamante Reconvergente Controlada (RDC)}
\author{Projeto RDC}
\date{\today}

\begin{document}
	\maketitle
	
	\section{Introdução}
	
	Este documento define a \textbf{API mínima obrigatória} da Estrutura de Dados
	Diamante Reconvergente Controlada (RDC). A API aqui descrita estabelece os
	\textbf{contratos formais} entre arquitetura e implementação.
	
	Nenhuma função adicional pode violar as invariantes definidas em
	\texttt{SPEC\_RDC.md} e \texttt{INVARIANTS\_RDC.tex}.
	
	\section{Princípios da API}
	
	\begin{itemize}[leftmargin=2em]
		\item A API deve ser mínima e suficiente
		\item Todas as operações devem ser determinísticas
		\item Nenhuma função pode expor estados internos inconsistentes
		\item A API deve permitir testes de propriedades
	\end{itemize}
	
	\section{Tipos Fundamentais}
	
	\begin{itemize}[leftmargin=2em]
		\item \texttt{Core}
		\item \texttt{Split}
		\item \texttt{NodoTransistor}
		\item \texttt{Merge}
		\item \texttt{DiamanteRDC}
	\end{itemize}
	
	\section{Funções da API}
	
	\subsection{create\_core}
	
	\textbf{Assinatura conceitual:}
	
	\begin{verbatim}
		Core create_core(Data payload)
	\end{verbatim}
	
	\textbf{Contrato:}
	\begin{itemize}
		\item Cria exatamente um Core
		\item O Core é imutável após a criação
	\end{itemize}
	
	\subsection{split}
	
	\textbf{Assinatura conceitual:}
	
	\begin{verbatim}
		Split split(Core core)
	\end{verbatim}
	
	\textbf{Contrato:}
	\begin{itemize}
		\item Deriva visões do Core
		\item Não modifica o Core
	\end{itemize}
	
	\subsection{evaluate\_node}
	
	\textbf{Assinatura conceitual:}
	
	\begin{verbatim}
		State evaluate_node(NodoTransistor node)
	\end{verbatim}
	
	\textbf{Contrato:}
	\begin{itemize}
		\item Executa o gate matemático
		\item Atualiza o estado do nó
		\item Não produz efeitos colaterais externos
	\end{itemize}
	
	\subsection{step}
	
	\textbf{Assinatura conceitual:}
	
	\begin{verbatim}
		void step(DiamanteRDC d)
	\end{verbatim}
	
	\textbf{Contrato:}
	\begin{itemize}
		\item Executa um ciclo completo da RDC
		\item Avalia nós ativos
		\item Não viola invariantes
	\end{itemize}
	
	\subsection{merge}
	
	\textbf{Assinatura conceitual:}
	
	\begin{verbatim}
		Result merge(DiamanteRDC d)
	\end{verbatim}
	
	\textbf{Contrato:}
	\begin{itemize}
		\item Realiza reconvergência determinística
		\item Independe da ordem de avaliação
	\end{itemize}
	
	\subsection{collapse\_node}
	
	\textbf{Assinatura conceitual:}
	
	\begin{verbatim}
		void collapse_node(NodoTransistor node)
	\end{verbatim}
	
	\textbf{Contrato:}
	\begin{itemize}
		\item Reduz o nó ao estado COLLAPSE
		\item Libera recursos
	\end{itemize}
	
	\subsection{reset}
	
	\textbf{Assinatura conceitual:}
	
	\begin{verbatim}
		void reset(DiamanteRDC d)
	\end{verbatim}
	
	\textbf{Contrato:}
	\begin{itemize}
		\item Retorna a RDC ao estado inicial
		\item Não altera a SPEC
	\end{itemize}
	
	\section{Garantias da API}
	
	A API garante:
	
	\begin{itemize}[leftmargin=2em]
		\item Determinismo
		\item Reconvergência obrigatória
		\item Controle matemático local
		\item Previsibilidade de memória
	\end{itemize}
	
	\section{Conclusão}
	
	Esta API constitui o \textbf{limite formal} entre arquitetura e implementação.
	Qualquer código que respeite estes contratos e as invariantes associadas pode ser
	considerado uma implementação válida da RDC.
	
\end{document}
